% !TITLE 天净沙·秋思
% !AUTHOR 马致远
% !DYNASTY 元
% !GENRE 散曲
% !ORDER 1

\begin{poem}
枯藤\textsuperscript{1}老树昏鸦\textsuperscript{2}，小桥流水人家，古道\textsuperscript{3}西风瘦马。
夕阳西下，断肠人\textsuperscript{4}在天涯\textsuperscript{5}。
\end{poem}

\begin{annotations}
\notetext{1}{枯藤：干枯的藤蔓。枯萎的植物暗示深秋的萧瑟。}
\notetext{2}{昏鸦：黄昏时归巢的乌鸦。昏，指天色昏暗。乌鸦在古诗词中常代表荒凉和不祥。}
\notetext{3}{古道：古老荒凉的道路。与“西风”“瘦马”组合，写旅途的艰辛。}
\notetext{4}{断肠人：形容极度悲伤的人。断肠，肝肠寸断，古人认为悲伤至极会使肠子断裂。}
\notetext{5}{天涯：天边，极远的地方。在天涯，意味着远离家乡亲人。}
\end{annotations}

\begin{appreciation}
《天净沙·秋思》被后人称为"秋思之祖"，是元曲小令里最有名的一首之一。二十八个字，写尽一个旅人的黄昏。

"枯藤老树昏鸦，小桥流水人家，古道西风瘦马"——三句话，九个名词，没有一个动词。枯藤、老树、黄昏的乌鸦，是萧瑟；小桥、流水、住着人的房子，是温暖；古道、西风、瘦马，是赶路。九个词连起来，像镜头慢慢扫过一圈：先扫过荒凉，再扫过别人家的灯火，最后停在这匹瘦马和马上的人身上。一句"小桥流水人家"是全曲最狠的一笔——不是写他看见了什么，而是写他看见了别人的家。乌鸦知道回巢，人家知道生火做饭，而他还在这条古道上：风是西风，马是瘦马，人是断肠人。

"夕阳西下"把时间钉死在黄昏，然后"断肠人在天涯"——肝肠寸断的人，在天边。前面堆了那么多景物，一个字不写心情，到最后一句才把人放出来。原来枯藤、老树、昏鸦、瘦马、西风，全是替这个人在难受。

还记得《蒹葭》里的追寻者吗？"溯洄从之，道阻且长"，他还在路上。还记得《行行重行行》里的游子吗？"相去万余里，各在天一涯"。中国文学里一直有一个在路上的人：汉代的游子走到"衣带日以缓"，元代的旅人走到"断肠人在天涯"，走了一千年，家还是那么远。马致远自己早年也漂泊，后来取号"东篱"，算是安了家，可这首小令里的黄昏，他替所有赶路的人留了一辈子。

白朴的《天净沙·秋》写秋天好看的一面，这一首写秋天难熬的一面。同一个曲牌，同是黄昏，同有寒鸦，一个在风景里，一个在路上。

哥哥希望你永远不用懂"断肠人在天涯"这句话。但万一将来你也在外地，夜深了，路灯下一个人往回走——记得天冷了加件衣服，想家了就给家里打个电话。天涯再远，电话线那头总有人在。
\end{appreciation}
